\documentclass[11pt,reqno]{amsart}
\usepackage[T1]{fontenc}
\usepackage[utf8]{inputenc}
\usepackage{lmodern,amssymb,booktabs,array,hyperref,xurl}
\usepackage[a4paper,margin=27mm]{geometry}
\hypersetup{colorlinks=true,linkcolor=black,citecolor=black,urlcolor=blue}
\emergencystretch=1em
\newtheorem{maintheorem}{Theorem}
\newtheorem{maincorollary}[maintheorem]{Corollary}
\newtheorem{lemma}{Lemma}[section]
\newtheorem{corollary}[lemma]{Corollary}
\newtheorem{proposition}[lemma]{Proposition}
\title[Antipodal necessity through degree three]{Rotation symmetric bent functions of degree at most three: the antipodal term is necessary}
\author{Filipe Martins Vono}
\subjclass[2020]{Primary 94D10; Secondary 06E30, 11T71}
\keywords{Bent functions, rotation symmetric Boolean functions, homogeneous Boolean functions, algebraic normal form, quadratic forms in characteristic two}
\date{Editorial revision, October 19, 2026}
\begin{document}
\begin{abstract}
A Boolean function on $\mathbb F_2^n$ is rotation symmetric if it is invariant under cyclic shifts of its coordinates. St\u{a}nic\u{a} and Maitra conjectured that no homogeneous rotation symmetric function of degree at least three is bent. We prove the cubic case in every even dimension. More precisely, if $n>0$ is even and $f$ is a rotation symmetric bent function of degree at most three, then the coefficient of $x_0x_{n/2}$ in the algebraic normal form of $f$ equals $1$. A homogeneous cubic has this coefficient equal to zero, so it cannot be bent. The proof has three ingredients: a reduction to the fixed space of an odd-order rotation, obtained by averaging on the radical of a quadratic form; a rigidity statement for rotation symmetric quadratics in dimension a power of two; and a contradiction between the balance and the constancy of one fibre of the function. No case analysis in $n$ is needed.
\end{abstract}
\maketitle
\section{Introduction}

A Boolean function $f:\mathbb F_2^n\to\mathbb F_2$, $n$ even, is \emph{bent} if all its Walsh coefficients have the same absolute value $2^{n/2}$. Bent functions, introduced by Rothaus \cite{ref6}, attain the largest possible nonlinearity in even dimension, where nonlinearity is the minimum Hamming distance to an affine Boolean function.

A function is \emph{rotation symmetric} if $f(x_0,\dots,x_{n-1})=f(x_1,\dots,x_{n-1},x_0)$. Their truth tables are constant on cyclic orbits, which reduces the number of independent values and is useful in computational searches. Rotation symmetric bent functions (RSBFs) have been studied extensively. A rotation symmetric function is \emph{homogeneous} of degree $d$ if its algebraic normal form (ANF) contains only monomials of degree $d$. Quadratic homogeneous RSBFs exist in every even dimension. St\u{a}nic\u{a} and Maitra \cite{ref1} conjectured that no homogeneous RSBF of degree $\ge 3$ exists. Earlier work includes Meng, Chen and Fu \cite{ref2}, Zhang and Gao \cite{ref3}, Cusick and Sanger \cite{ref4}, and Sun, Shi, Liu and Fu \cite{ref5}. In particular, the homogeneous cubic case in dimensions $n=2p$, with $p>2$ prime, follows from \cite[Theorem~3.8]{ref4}. The present statement covers every positive even dimension and allows lower-degree terms. Cusick and Sanger also point out that Theorem 3.1 of \cite{ref3} is false; that criterion is not used here.

\textbf{Main result.} For even $n=2H$ let $[x_0x_H]f\in\mathbb F_2$ denote the coefficient of the monomial $x_0x_H$ in the ANF of $f$. For rotation symmetric $f$ this equals the coefficient of every \emph{antipodal} monomial $x_ix_{i+H}$.

\begin{maintheorem}Let $n>0$ be even and let $f:\mathbb F_2^n\to\mathbb F_2$ be rotation symmetric, bent, and of degree at most $3$. Then $[x_0x_{n/2}]f=1$.\end{maintheorem}

\begin{maincorollary}For every even $n>0$ there is no homogeneous rotation symmetric bent function of degree $3$.\end{maincorollary}

The corollary is immediate: a homogeneous cubic has no quadratic monomials. Theorem 1 is stronger than the corollary because it allows arbitrary lower-degree terms, and its hypotheses are nonvacuous, as Section 5 shows: the antipodal form $\sum_{i<H}x_ix_{i+H}$ plus any affine rotation symmetric term is bent.

\textbf{Idea of the proof.} Write $n=2^s m$ with $m$ odd. (i) Bentness of a cubic rotation symmetric $f$ passes to its restriction $g$ to the subspace of vectors fixed by rotation through $2^s$ coordinates, a rotation symmetric function on $\mathbb F_2^{2^s}$ (Section 3). This uses that a quadratic is balanced exactly when its normalisation $q+q(0)$ is non-zero on the radical of its polar, and that averaging over a group action of odd order does not change this behaviour. (ii) If $[x_0x_H]f=0$, terms of the ANF cancel in pairs under the half-turn $x_i\mapsto x_{i+H}$, so $g$ is constant on the diagonal $\{(u,u)\}$ (Section 5). (iii) On $\mathbb F_2^{2^s}$, a bent function of degree $\le 3$ that is constant on the diagonal cannot exist: Parseval forces a certain fibre to be balanced, while a rigidity property of rotation symmetric quadratics forces the same fibre to be constant (Section 4). The power-of-two assumption enters the quadratic rigidity argument through nilpotence of $\rho+I$. Cyclic invariance is also used in the fixed-space coordinates, half-turn cancellation and fibre covariance.

\textbf{Machine verification.} The statement of Theorem 1 has been checked in Lean 4; see Section 6 for exactly what was and was not formalized. The paper is self-contained and the human proof below does not depend on the formalization.

\textbf{Organization.} Section 2 fixes notation and collects standard facts. Section 3 treats quadratic forms and the odd-order reduction. Section 4 proves the obstruction in power-of-two dimension. Section 5 proves Theorem 1 and discusses nonvacuity and scope. Section 6 describes the formal verification.


\section{Preliminaries}

\subsection{Notation}

All spaces are finite-dimensional over $\mathbb F_2$. A function has \emph{degree at most $d$} if its ANF has degree $\le d$; we call such functions \emph{quadratic} ($d=2$) or \emph{cubic} ($d=3$) without implying that the top degree is attained. Composing with a linear map does not raise degree.

\begin{center}\small\begin{tabular}{@{}>{\raggedright\arraybackslash}p{0.29\textwidth}>{\raggedright\arraybackslash}p{0.65\textwidth}@{}}\toprule
Symbol & Meaning \\\midrule
$V_n=\mathbb F_2^n$ & indices of coordinates are read modulo $n$ \\
$\rho$ & rotation, $(\rho x)_i=x_{i+1}$; $f$ rotation symmetric means $f\circ\rho=f$ \\
$c_A(f)$ & ANF coefficient of $\prod_{i\in A}x_i$ \\
$D_af(x)=f(x+a)+f(x)$ & derivative \\
$W_f(b)=\sum_x(-1)^{f(x)+b\cdot x}$ & Walsh transform \\
$\Delta_f(a)=\sum_x(-1)^{D_af(x)}$ & autocorrelation \\
$G(q)=\sum_x(-1)^{q(x)}$ & character sum \\
$B_q(a,b)=q(a+b)+q(a)+q(b)+q(0)$ & polar form \\
$\operatorname{Rad}B=\{r: B(r,\cdot)=0\}$ & radical \\
$q_0=q+q(0)$ & normalised quadratic \\
\bottomrule\end{tabular}\end{center}

A function $f$ on $V_n$ is \textbf{bent} if $n$ is even and $|W_f(b)|=2^{n/2}$ for all $b$. The same definition is used on any space $U\cong V_k$ after choosing a basis; it does not depend on the basis.

\subsection{Standard facts}

\begin{lemma}[derivatives]

(a) If $\deg q\le 2$, then $B_q$ is alternating bilinear and $B_q(a,b)=D_aD_bq(x)$ for every $x$.
(b) If $\deg g\le3$, then $\Theta(a,b,c)=D_aD_bD_cg(x)$ does not depend on $x$ and is symmetric and trilinear. Moreover $D_aD_bg(x)=D_aD_bg(0)+\Theta(a,b,x)$, and $B_{D_cg}(a,b)=\Theta(a,b,c)$.

\end{lemma}
\begin{proof}
A derivative lowers the degree by at least one. Thus $D_aD_bq$ and $D_aD_bD_cg$ are constant, giving the asserted independence of the base point; the formula $D_aD_bq(0)=B_q(a,b)$ is a rewriting of the definition. The identity $D_{a+a'}=D_a+D_{a'}+D_{a'}D_a$ shows additivity in each argument, since the correction term is one more derivative of the constant function and vanishes. Alternation follows from $B_q(a,a)=q(0)+2q(a)+q(0)=0$. Symmetry holds because derivatives commute. The remaining formula is $D_xD_aD_bg(0)=D_aD_bg(x)+D_aD_bg(0)$.
\end{proof}

\begin{lemma}[flat Walsh spectrum $\Leftrightarrow$ balanced derivatives]
For any $f$ on $V_n$ and any $b$,

\begin{equation*}W_f(b)^2=\sum_{a}\Delta_f(a)(-1)^{b\cdot a}. \tag{2.1}\end{equation*}

Consequently $W_f(b)^2=2^n$ for all $b$ if and only if $\Delta_f(a)=0$ for every $a\ne0$, that is, if and only if every nonzero derivative $D_af$ is balanced.

\end{lemma}
\begin{proof}
Expand $W_f(b)^2=\sum_{x,y}(-1)^{f(x)+f(y)+b\cdot(x+y)}$ and put $y=x+a$. This is the classical derivative criterion for bentness \cite{ref6}. For the equivalence, Fourier inversion gives $\Delta_f(a)=2^{-n}\sum_bW_f(b)^2(-1)^{a\cdot b}$. If $W_f^2\equiv2^n$, the sum vanishes for $a\neq0$. Conversely, if $\Delta_f(a)=0$ for $a\ne0$ then (2.1) gives $W_f(b)^2=\Delta_f(0)=2^n$.
\end{proof}

\begin{lemma}[rotation symmetry and the ANF]
If $f\circ\rho=f$ then $c_{A+1}(f)=c_A(f)$ for every $A\subseteq\mathbb Z_n$. In particular, for $n=2H$, all coefficients $c_{\{i,i+H\}}(f)$ are equal to $\alpha(f):=[x_0x_H]f$.

\end{lemma}
\begin{proof}
$(f\circ\rho)(x)=\sum_Ac_A(f)\prod_{i\in A}x_{i+1}$ is an ANF, and ANFs are unique.
\end{proof}

\emph{Remark.} $\alpha(f)$ is a single coefficient. It should not be confused with the sum of all antipodal coefficients, which in characteristic two equals $H\,\alpha(f)$ and vanishes for even $H$.


\section{Quadratic forms and the odd-order reduction}

\begin{lemma}[Gauss sums]
Let $q$ be a function of degree $\le2$ on a space $V$, and $R=\operatorname{Rad}B_q$. Then $q_0|_R$ is linear and

\begin{equation*}G(q)^2=|V|\sum_{r\in R}(-1)^{q_0(r)}. \tag{3.1}\end{equation*}

Consequently

\begin{equation*}q\ \text{is balanced}\iff q_0|_R\not\equiv0. \tag{3.2}\end{equation*}


\end{lemma}
\begin{proof}
For $r,r'\in R$, $q_0(r+r')=q_0(r)+q_0(r')+B_q(r,r')=q_0(r)+q_0(r')$, so $q_0|_R$ is linear. Write $G(q)^2=\sum_{x,y}(-1)^{q(x)+q(y)}$ and substitute $y=x+r$. Since $q(x+r)+q(x)=B_q(x,r)+q_0(r)$,

\begin{equation*}G(q)^2=\sum_{r\in V}(-1)^{q_0(r)}\sum_{x}(-1)^{B_q(x,r)}.\end{equation*}

The inner sum is $|V|$ if $r\in R$ and $0$ otherwise, which is (3.1). The sum of a nontrivial character over $R$ vanishes and the sum of the trivial one equals $|R|>0$; so $G(q)=0$ iff $q_0|_R\ne0$.
\end{proof}

(The cases of constant $q$, degenerate $B_q$ and $\dim V=0$ are included.)

\begin{lemma}[odd-order averaging]
Let $T\in\mathrm{GL}(V)$ with $T^m=I$ for an odd integer $m\ge1$, let $U=\operatorname{Fix}(T)$, and let $q$ have degree $\le2$ with $q\circ T=q$. Then $q$ is balanced on $V$ if and only if $q|_U$ is balanced.

\end{lemma}
\begin{proof}
Put $\Pi=I+T+\dots+T^{m-1}$. Since $T\Pi=\Pi$, $\Pi(V)\subseteq U$, and $\Pi|_U=mI=I$ because $m$ is odd; hence $\Pi(V)=U$. The invariance $q\circ T=q$ gives $B_q(Ta,Tb)=B_q(a,b)$, so $R=\operatorname{Rad}B_q$ is $T$-stable.

\emph{Radical of the restriction.} For $u\in U$ and $x\in V$,

\begin{equation*}B_q(u,\Pi x)=\sum_{j<m}B_q(T^{-j}u,x)=mB_q(u,x)=B_q(u,x).\end{equation*}

If $u\in U$ is orthogonal to $U$, it is orthogonal to every $\Pi x$, hence to every $x$. So $\operatorname{Rad}(B_q|_{U\times U})=R\cap U$.

\emph{Transfer of the nonvanishing.} The normalised function $(q|_U)_0$ is the restriction of $q_0$. If $q_0$ is nonzero on $R\cap U$ it is nonzero on $R$. Conversely let $r\in R$ with $q_0(r)=1$. Then $\Pi r\in R\cap U$, and by linearity of $q_0|_R$ and the invariance $q_0\circ T=q_0$,

\begin{equation*}q_0(\Pi r)=\sum_{j<m}q_0(T^jr)=m\,q_0(r)=1.\end{equation*}

Thus $q_0|_R\ne0$ iff $q_0|_{R\cap U}\ne0$, and (3.2) on $V$ and on $U$ concludes.
\end{proof}

\begin{corollary}[bent restriction]
Let $n$ be even, $T$ as in Lemma 3.2 on $V=V_n$, and $f:V_n\to\mathbb F_2$ bent with $\deg f\le3$ and $f\circ T=f$. Then $\dim U$ is even and $f|_U$ is bent.

\end{corollary}
\begin{proof}
Let $0\ne a\in U$. Since $Ta=a$, $D_af\circ T=D_a(f\circ T)=D_af$, and $D_af$ has degree $\le2$ and is balanced (Lemma 2.2). By Lemma 3.2 its restriction $D_a(f|_U)=(D_af)|_U$ is balanced. So $\Delta_{f|_U}(a)=0$ for every nonzero $a\in U$, and Lemma 2.2 (applied on $U\cong V_k$, $k=\dim U$) gives $W_{f|_U}(b)^2=2^k$ for all $b$. Because Walsh coefficients are integers, $k$ is even, and $f|_U$ is bent.
\end{proof}

\begin{lemma}[rigidity of rotation symmetric quadratics]
Let $N=2^s$, $s\ge0$. If $q:V_N\to\mathbb F_2$ has degree $\le2$, satisfies $q\circ\rho=q$ and is balanced, then $B_q=0$.

\end{lemma}
\begin{proof}
Suppose $B:=B_q\ne0$; then $R=\operatorname{Rad}B$ is a proper subspace, stable under $\rho$ (as $B$ is $\rho$-invariant). Let $\nu=\rho+I$. In characteristic two with $N$ a power of two, $\nu^N=\rho^N+I=0$. Moreover

\begin{equation*}\operatorname{im}\nu=\Big\{x:\ \textstyle\sum_ix_i=0\Big\}:\end{equation*}

the inclusion $\subseteq$ is clear, and given $x$ of even weight the equations $y_{i+1}+y_i=x_i$ can be solved successively around the cycle, the wrap-around condition being exactly $\sum x_i=0$.

\emph{Step 1: $R\subseteq\operatorname{im}\nu$.} If $R$ contained a vector of odd weight, then since $\operatorname{im}\nu$ is a hyperplane not containing it, $V_N=R+\nu V_N$. Using $\nu R\subseteq R$ we get $\nu V_N=\nu R+\nu^2V_N\subseteq R+\nu^2V_N$, so $V_N=R+\nu^2V_N$, and by induction $V_N=R+\nu^jV_N$ for all $j$. With $j=N$ this gives $V_N=R$, a contradiction.

\emph{Step 2: $q_0|_R=0$.} Let $r\in R$, so $r=\rho y+y$ for some $y$ by Step 1. Using $q_0(\rho y)=q_0(y)$ and the alternation of $B$,

\begin{equation*}q_0(r)=q_0(\rho y)+q_0(y)+B(\rho y,y)=B(r+y,y)=B(r,y)+B(y,y)=0.\end{equation*}

By (3.2) $q$ is not balanced, a contradiction.
\end{proof}


\section{The constant-diagonal obstruction}

Throughout this section $N=2h$ and a vector of $V_N$ is written $(u,w)$ with $u,w\in V_h$. Let $\mathbf 1_h$ denote the all-ones vector of $V_h$, $J=\mathbf 1_N$, and put

\begin{equation*}d(u)=(u,u),\qquad \varepsilon(w)=(0,w).\end{equation*}

Let $\theta(u,w)=(w,u)$ be the block swap; it equals $\rho^h$.

\begin{lemma}[fibres of a bent function with constant diagonal]
Let $g:V_N\to\mathbb F_2$ be bent and constant on $\{d(u)\}$. For $z\in V_h$ let $\varphi_z(u)=g(u,u+z)$. Then $\varphi_z$ is balanced for every $z\ne0$.

\end{lemma}
\begin{proof}
Put $\gamma(z)=\sum_u(-1)^{\varphi_z(u)}$. Its Walsh transform in $z$ is

\begin{equation*}\sum_{z}\gamma(z)(-1)^{b\cdot z}=\sum_{u,v}(-1)^{g(u,v)+b\cdot(u+v)}=W_g(b,b),\end{equation*}

where $v=u+z$. Parseval on $V_h$ and bentness ($W_g^2=2^{N}=2^{2h}$) give

\begin{equation*}\sum_z\gamma(z)^2=2^{-h}\sum_{b\in V_h}W_g(b,b)^2=2^{-h}\cdot2^h\cdot2^{2h}=2^{2h}.\end{equation*}

Since $g$ is constant on the diagonal, $\gamma(0)^2=(2^h)^2=2^{2h}$. Hence $\gamma(z)=0$ for all $z\ne0$.
\end{proof}

The hypothesis of constancy on the diagonal is indispensable: without it $\gamma(0)^2$ may be smaller and the argument gives nothing.

\begin{proposition}
Let $N=2^s$ with $s\ge1$. A rotation symmetric function $g:V_N\to\mathbb F_2$ of degree at most three that is constant on the diagonal is not bent.

\end{proposition}
\begin{proof}
Assume $g$ is bent. We show that the fibre $\varphi:=\varphi_{\mathbf 1_h}$ is constant, which contradicts Lemma 4.1 because $\mathbf 1_h\neq0$.

\emph{(i) The derivative in the direction $J$.} $D_Jg$ has degree $\le2$, is $\rho$-invariant ($\rho J=J$) and is balanced since $g$ is bent. By Lemma 3.4, $B_{D_Jg}=0$, that is, by Lemma 2.1(b),

\begin{equation*}\Theta(a,b,J)=0\quad\text{for all }a,b. \tag{4.1}\end{equation*}


\emph{(ii) The polar of $\varphi$.} Since $(u,u+\mathbf 1_h)=d(u)+\varepsilon(\mathbf 1_h)$,

\begin{equation*}B_\varphi(u,v)=D_{d(u)}D_{d(v)}g\big(\varepsilon(\mathbf 1_h)\big)=D_{d(u)}D_{d(v)}g(0)+\Theta\big(d(u),d(v),\varepsilon(\mathbf 1_h)\big).\end{equation*}

The first term is $g(d(u+v))+g(d(u))+g(d(v))+g(0)=0$ because $g$ is constant on the diagonal. Hence

\begin{equation*}B_\varphi(u,v)=\Theta\big(d(u),d(v),\varepsilon(\mathbf 1_h)\big). \tag{4.2}\end{equation*}


\emph{(iii) Use of the block swap.} $\theta=\rho^h$ preserves $g$, hence $\Theta$: $\Theta(\theta a,\theta b,\theta c)=\Theta(a,b,c)$. It fixes $d(u)$ and satisfies

\begin{equation*}J=\varepsilon(\mathbf 1_h)+\theta\varepsilon(\mathbf 1_h),\qquad d(v)=\varepsilon(v)+\theta\varepsilon(v).\end{equation*}

Trilinearity and $\theta$-invariance give

\begin{equation*}\begin{aligned}\Theta\big(d(u),\varepsilon(v),J\big)&=\Theta\big(d(u),\varepsilon(v),\varepsilon(\mathbf 1_h)\big)\\&\quad+\Theta\big(d(u),\theta\varepsilon(v),\varepsilon(\mathbf 1_h)\big)\\&=\Theta\big(d(u),d(v),\varepsilon(\mathbf 1_h)\big),\end{aligned}\end{equation*}

where the second summand was obtained by applying $\theta$ to all three arguments. With (4.2) and (4.1),

\begin{equation*}B_\varphi(u,v)=\Theta\big(d(u),\varepsilon(v),J\big)=0.\end{equation*}


\emph{(iv) Constancy of $\varphi$.} Since $B_\varphi=0$, $\varphi+\varphi(0)$ is additive, hence linear: $\varphi(u)=\ell(u)+c$. Let $\sigma$ be the rotation of $V_h$ and $e$ its last standard basis vector $e_{h-1}$. A direct computation gives $\rho(u,u+\mathbf 1_h)=(\sigma u+e,\ \sigma u+e+\mathbf 1_h)$, so $\varphi(\sigma u+e)=\varphi(u)$, i.e. $\ell(\sigma u)+\ell(e)=\ell(u)$ for all $u$. Writing $\ell(u)=\sum\lambda_iu_i$, comparing coefficients gives $\lambda_{i-1}=\lambda_i$ for all $i$ (so all $\lambda_i$ are equal to a common $\lambda$), and $u=0$ gives $\ell(e)=\lambda=0$. Thus $\ell=0$ and $\varphi$ is constant. (For $h=1$, $\sigma=I$ and the same computation gives $\ell(e)=0$.)
\end{proof}


\section{Proof of the main theorem}

\begin{lemma}[half-turn cancellation]
Let $n=2H$ and let $f$ be rotation symmetric of degree at most $3$. Then for every $u\in V_H$,

\begin{equation*}f(u,u)=f(0)+\alpha(f)\sum_{i<H}u_i. \tag{5.1}\end{equation*}


\end{lemma}
\begin{proof}
Let $\tau=\rho^H$ be the half-turn, and $x=(u,u)$, so $\tau x=x$ and $x^A:=\prod_{i\in A}x_i$ equals $x^{\tau A}$. By Lemma 2.3, $c_A(f)=c_{\tau A}(f)$. For $A\ne\tau A$ the monomials of $A$ and $\tau A$ contribute $c_Ax^A+c_{\tau A}x^{\tau A}=0$. The sets with $\tau A=A$ are unions of antipodal pairs $\{i,i+H\}$; having size $\le3$, they are $\emptyset$ or a single pair. The empty set contributes $f(0)$ and the pair $\{i,i+H\}$ contributes $c_{\{i,i+H\}}(f)\,x_ix_{i+H}=\alpha(f)u_i$ (as $u_i^2=u_i$).
\end{proof}

No assumption that the cyclic orbits of $A$ have full length is made: the argument only uses the pairing $A\leftrightarrow\tau A$ and the classification of $\tau$-fixed sets of size $\le3$.

\begin{proof}[Proof of Theorem 1]
Write $n=Nm$ with $N=2^s$, $s\ge1$, and $m$ odd, and put $N=2h$. Let $T=\rho^N$; then $T^m=\rho^n=I$, and $U=\operatorname{Fix}(T)=\{x:x_{i+N}=x_i\ \forall i\}$. The map

\begin{equation*}\iota:V_N\to U,\qquad(\iota y)_i=y_{i\bmod N}\end{equation*}

is a linear isomorphism that intertwines the rotations of $V_N$ and $V_n$. Hence $g:=f\circ\iota$ is rotation symmetric of degree $\le3$, and bent by Corollary 3.3 (bentness is invariant under linear isomorphisms).

Assume $\alpha(f)=0$. By (5.1), $f$ is constant, equal to $f(0)$, on the vectors $(v,v)$, $v\in V_{n/2}$, i.e. on the vectors fixed by the half-turn $\rho^{n/2}$. For $u\in V_h$ the vector $\iota(d(u))$ has period dividing $h$ and $h\mid n/2=hm$, so it is fixed by the half-turn. Therefore

\begin{equation*}g(d(u))=f(0)\qquad\text{for all }u\in V_h,\end{equation*}

and $g$ is constant on the diagonal. This contradicts Proposition 4.2. Hence $\alpha(f)=1$.
\end{proof}

Corollary 2 follows because a homogeneous cubic has $c_A=0$ for $|A|=2$.

\textbf{Nonvacuity.} For every $H\ge1$ and $a,c\in\mathbb F_2$ the function

\begin{equation*}\kappa_H(x)+a\sum_ix_i+c,\qquad\kappa_H(x)=\sum_{i<H}x_ix_{i+H},\end{equation*}

is rotation symmetric and bent: its Walsh transform factors into $H$ independent two-variable transforms $\sum_{s,t}(-1)^{st+b_1s+b_2t}=\pm2$, and affine additions permute Walsh frequencies and may reverse signs. Thus the hypotheses of Theorem 1 are attained, including invariant affine terms. The theorem gives a \emph{necessary} coefficient condition, not a classification of bent cubics.

\textbf{Example $n=6$.} Here $N=2$, $m=3$, $U=\{(a,b,a,b,a,b)\}$ and $g(a,b)=f(a,b,a,b,a,b)$ is a symmetric function of two variables, hence $g=\beta ab+\gamma(a+b)+\delta$. It is bent iff $\beta=1$, and $g(0,0)+g(1,1)=\beta$. If $\alpha(f)=0$, both $(0,0,0,0,0,0)$ and $(1,1,1,1,1,1)$ are half-turn fixed, so $g(0,0)=g(1,1)=f(0)$; then $\beta=0$ and $g$ is affine, hence not bent. In general $N>2$ the same mechanism is carried by Proposition 4.2.

\textbf{Scope.} The proof is sensitive to the degree bound in two places. For a quartic, derivatives are cubic and Lemma 3.2 does not apply directly; and a $\tau$-fixed support may contain two antipodal pairs, so (5.1) changes. No statement about degree $\ge4$ follows.


\section{Formal verification}

The statement of Theorem 1 has been formalized in Lean 4 (toolchain \nolinkurl{leanprover/lean4:v4.15.0}, Mathlib commit \nolinkurl{9837ca9d65d9de6fad1ef4381750ca688774e608}) in the repository \nolinkurl{hrsbf-bent-verification}, commit \nolinkurl{6dd183685ef8a5b033c1bcf8733496d57b853cc7}, theorem \nolinkurl{HRSBF.Antipodal.universal_antipodal_walsh_absolute}. Its hypotheses are: $H>0$, $f:V_{2H}\to\mathbb F_2$ with canonical ANF of degree $\le3$, rotation symmetry, and $|W_f(b)|=2^H$ for all $b$; its conclusion is that the ANF coefficient of $\{0,H\}$ equals $1$. The development depends only on the axioms \nolinkurl{propext}, \nolinkurl{Classical.choice}, \nolinkurl{Quot.sound}.

The formal proof of the odd-order reduction decomposes $V$ into $\operatorname{Fix}(T)$ and $\ker\Pi$, which differs from the radical-averaging proof of Lemma 3.2 given here. The latter has \emph{not} been separately formalized; its complete mathematical argument is given in Lemma 3.2. The formalization does not use a general formula for the behaviour of ANF coefficients under repetition of blocks, and neither does the argument above. Finite computations (Gauss identity for 2196 quadratics, invariants of an order-three action, equivalence of vanishing polar and affinity) were run as sanity checks; they play no role in the proof.

\begin{thebibliography}{6}
\raggedright
\bibitem{ref1} P.~St\u{a}nic\u{a} and S.~Maitra, Rotation symmetric Boolean functions---count and cryptographic properties, \emph{Discrete Applied Mathematics} \textbf{156} (2008), no.~10, 1567--1580. \href{https://doi.org/10.1016/j.dam.2007.04.029}{doi:10.1016/j.dam.2007.04.029}.
\bibitem{ref2} Q.~Meng, L.~Chen and F.-W.~Fu, On homogeneous rotation symmetric bent functions, \emph{Discrete Applied Mathematics} \textbf{158} (2010), no.~10, 1111--1117. \href{https://doi.org/10.1016/j.dam.2010.02.009}{doi:10.1016/j.dam.2010.02.009}.
\bibitem{ref3} X.~Zhang and G.~Gao, On the conjecture about the nonexistence of rotation symmetric bent functions, preprint (2013), \href{https://arxiv.org/abs/1303.2282}{arXiv:1303.2282}.
\bibitem{ref4} T.~W.~Cusick and E.~M.~Sanger, Rotation symmetric bent Boolean functions for $n=2p$, preprint (2017), \href{https://arxiv.org/abs/1708.09313}{arXiv:1708.09313}.
\bibitem{ref5} L.~Sun, Z.~Shi, J.~Liu and F.-W.~Fu, On the conjecture about the nonexistence of homogeneous rotation symmetric bent functions, \emph{Designs, Codes and Cryptography} \textbf{94} (2026), no.~4, article~93. \href{https://doi.org/10.1007/s10623-026-01848-4}{doi:10.1007/s10623-026-01848-4}.
\bibitem{ref6} O.~S.~Rothaus, On ``bent'' functions, \emph{Journal of Combinatorial Theory, Series A} \textbf{20} (1976), no.~3, 300--305. \href{https://doi.org/10.1016/0097-3165(76)90024-8}{doi:10.1016/0097-3165(76)90024-8}.
\end{thebibliography}
\end{document}
